\section{Sequences}

\begin{definition}[Sequence]
A sequence is a function $a : \N \to \R$, written $\{a_n\}_{n \ge 1}$ or $(a_n)$.
We say $a_n \to L$ (or $\lim_{n\to\infty} a_n = L$) if for every $\eps > 0$ there exists $N$
such that $|a_n - L| < \eps$ for all $n \ge N$.
\end{definition}

\subsection{Limit laws and the function-limit bridge}
All algebraic limit laws (sum, product, quotient with nonzero denominator) carry over
from function limits. The most useful bridge:

\begin{theorem}
If $f : [1,\infty) \to \R$ is such that $a_n = f(n)$ and $\lim_{x\to\infty} f(x) = L$,
then $\lim_{n\to\infty} a_n = L$.
\end{theorem}

So you can apply L'H\^opital's rule to $f(x)$ to evaluate $\lim_{n\to\infty} a_n$.

\subsection{Monotone convergence}
\begin{theorem}[Monotone Convergence]
A monotone sequence converges iff it is bounded.
\end{theorem}

\subsection{Standard limits}
\begin{keybox}[Memorize]
\[
\lim_{n\to\infty} \frac{1}{n^p} = 0\ (p > 0), \quad
\lim_{n\to\infty} r^n = 0 \text{ if } |r| < 1, \quad
\lim_{n\to\infty} \sqrt[n]{n} = 1,
\]
\[
\lim_{n\to\infty} \sqrt[n]{c} = 1\ (c > 0), \quad
\lim_{n\to\infty} \Bigl(1 + \tfrac{x}{n}\Bigr)^n = e^x, \quad
\lim_{n\to\infty} \frac{n!}{n^n} = 0.
\]
\end{keybox}

\begin{example}
$\displaystyle a_n = \frac{n^2 + 3}{2n^2 - n} \;\to\; \tfrac{1}{2}$ (divide top and bottom by $n^2$).
\end{example}

\begin{example}[Squeeze]
$\displaystyle 0 \le \frac{\sin n}{n} \le \frac{1}{n} \;\Rightarrow\; \frac{\sin n}{n} \to 0.$
\end{example}

\begin{pitfallbox}
\begin{itemize}[leftmargin=*,nosep]
  \item Convergence of $\{a_n\}$ has nothing to do with convergence of the \emph{series} $\sum a_n$
        --- that is a different question (next section).
  \item $a_n \to 0$ does not imply $\sum a_n$ converges (e.g.\ harmonic series).
\end{itemize}
\end{pitfallbox}
